% =====================================================================
%  Cover letter -- standalone LaTeX, paste into Editorial Manager.
% =====================================================================
\documentclass[11pt]{letter}
\usepackage[T1]{fontenc}
\usepackage[utf8]{inputenc}
\usepackage[margin=1in]{geometry}
\usepackage[hidelinks]{hyperref}
\usepackage{parskip}

\signature{Md Ibrahim Khalil\\Sensa AS, Norway\\ibrahim@sensa.no\\On behalf of all authors}
\address{Md Ibrahim Khalil\\Sensa AS\\Oslo, Norway\\ibrahim@sensa.no}

\begin{document}
\begin{letter}{The Editor\\\emph{International Journal of Disaster Risk Reduction}\\Elsevier}

\opening{Dear Editor,}

We are pleased to submit our manuscript, \emph{FloodLite: A CPU-and-Browser Edge-Deployment Benchmark for Lightweight Flood Segmentation}, for consideration as a Research Article in the \emph{International Journal of Disaster Risk Reduction}.

Flood-segmentation models are almost always judged by accuracy on a benchmark. The responders who would use them run laptops, tablets and browsers, where latency, quantization behaviour and false alarms decide whether a model is usable. Apart from GPU timings on NVIDIA Jetson boards, the flood-segmentation literature does not report these quantities for commodity hardware. Our manuscript measures them for three lightweight UNet students and a UNet$+$EfficientNet-B0 teacher, on an Apple Silicon CPU and in browser WebAssembly, in FP32 and INT8.

Three findings make the work relevant to disaster-risk practice:

\textbf{1.~Flood-only training data produces models that cannot say ``no flood''.} The public Flood Semantic Segmentation Dataset (FSSD) contains no water-free images. Evaluated zero-shot on overlap-verified external data, every FSSD-trained model predicts water on 72--100\% of dry scenes, a failure invisible to every FSSD metric and disqualifying for unattended monitoring. Retraining on a combined public corpus that includes water-free images cuts student false-alarm rates to 3--9\%.

\textbf{2.~Distillation helps only with headroom, and it transmits the teacher's errors.} On saturated FSSD, no response-, feature- or combined-distillation configuration improves on the no-distillation baseline. On the less saturated combined corpus, combined distillation raises flood-only IoU significantly for all three students, but distilled students also false-alarm on the teacher's false-alarm scenes more than twice as often as undistilled ones. Teachers should be audited for the deployment-relevant failure before they are used for distillation.

\textbf{3.~INT8 post-training quantization is architecture-dependent.} Under one static-QDQ recipe, EfficientNet-Lite0 loses 0.043 IoU while MobileNetV3-Small (HardSwish) and MobileViT-XXS (attention) collapse by 0.45--0.46 IoU, and quantized models degrade further under domain shift. The recommended configuration, EfficientNet-Lite0-UNet INT8 on ONNX Runtime CPU, runs at 58~FPS on a 2023 laptop CPU, $2.2\times$ faster than the teacher on the same runtime and $3.83\times$ smaller; in the browser, where ONNX Runtime Web~1.17 has no INT8 convolution, FP32 students reach 6--7.5~FPS.

The work fits IJDRR's scope because its contribution is deployment evidence for disaster-response practice rather than a new network design: the quantities we report are the ones a humanitarian data-science team needs to decide whether, where and how a flood-segmentation model can be fielded, including the data it must be trained on to be trusted unattended.

\textbf{Reproducibility and openness.} All code, ONNX exports, per-fold checkpoints, and benchmarking harnesses are released under MIT licence at \url{https://github.com/NoorTasfia/flood-segmentation}. A Zenodo-archived snapshot (DOI to be minted at acceptance) will be linked in the camera-ready. The FSSD dataset is publicly available on Kaggle.

\textbf{Conflicts of interest.} None to declare. \textbf{Funding.} The work received no external funding. We confirm that the manuscript has not been published previously, is not under consideration elsewhere, and has been read and approved by all listed authors.

We would be grateful for your consideration of this submission and welcome any feedback from the editorial team and reviewers.

\closing{Sincerely,}

\end{letter}
\end{document}
